% ============================================================
%  CHAPTER 9 — RESULTS AND ANALYSIS
% ============================================================
\chapter{Results and Analysis}

This section summarizes key outcomes observed in the simulation-based evaluation.

\section{Model Outcomes}

\begin{table}[H]
\centering
\caption{Summary of Model Outcomes}
\label{tab:model_outcomes}
\begin{tabularx}{\textwidth}{@{}lXX@{}}
\toprule
\textbf{Component} & \textbf{Setup} & \textbf{Outcome} \\
\midrule
Latency predictor & Random Forest trained on RT-IoT2022-derived features & Low MAE with stable real-time inference \\
Anomaly detector & One-Class SVM novelty scoring on packet-level features & Consistent flagging of injected outliers with controlled false positives \\
Forecaster & BiLSTM sliding-window forecaster over recent latency history & Short-horizon trend estimates (t+10) for proactive monitoring \\
\bottomrule
\end{tabularx}
\end{table}

Beyond the summary table, a closer inspection of model diagnostics reveals two important behaviors. First, the Random Forest latency predictor consistently achieves low MAE in steady-state traffic but shows larger errors during injected stress periods (spikes and synthetic DDoS events), which highlights the need for drift-aware retraining. Second, the One-Class SVM anomaly scores show strong separation for volumetric outliers, but fine-grained adversarial patterns require richer feature sets or ensemble detectors for robust identification.

\section{Reinforcement Learning Behavior}

Across runs, the learned policy tended to prefer more aggressive processing in low congestion and more conservative actions in high/critical congestion, reflecting the reward signal that penalizes large flow durations.

\begin{table}[H]
\centering
\caption{Representative Q-Learning Policy Preference}
\label{tab:qtable}
\begin{tabular}{@{}lcccc@{}}
\toprule
\textbf{Action} & \textbf{Low} & \textbf{Medium} & \textbf{High} & \textbf{Critical} \\
\midrule
Accelerate ($0.5\times$) & Preferred & Sometimes & Rare & Rare \\
Maintain ($1.0\times$) & Sometimes & Preferred & Sometimes & Rare \\
Throttle ($2.0\times$) & Rare & Rare & Preferred & Preferred \\
\bottomrule
\end{tabular}
\end{table}

\section{Federated Learning Convergence}

In the federated prototype, the global model error decreased across communication rounds, indicating successful aggregation of distributed updates without sharing raw samples. However, convergence speed varied with client heterogeneity: clients with limited local data (simulated small-edge nodes) contributed noisy updates, suggesting that adaptive client weighting or robust aggregation (e.g., trimmed mean) could improve stability. Experiment logs show that communications rounds 3--7 produced the largest incremental gains, after which marginal improvements diminished under the current training schedule.

\section{Performance Benchmarks}

\begin{table}[H]
\centering
\caption{System Performance Benchmarks (Representative)}
\label{tab:benchmarks}
\begin{tabularx}{\textwidth}{@{}lXX@{}}
\toprule
\textbf{Component} & \textbf{Metric} & \textbf{Observation} \\
\midrule
Embedded inference & Per-packet prediction latency & Sub-millisecond class-model execution inside the JVM \\
RL decision & Action selection overhead & Negligible compared to packet processing \\
Desktop dashboards & UI responsiveness & Stable when using fixed-size sliding windows \\
Web dashboard & Load/interaction & Lightweight analysis over exported telemetry \\
\bottomrule
\end{tabularx}
\end{table}

\section{Quantitative Evaluation}

To complement the qualitative observations above, this section reports concrete numerical results obtained from a representative simulation run of 1,047 processed packets over the cleaned RT-IoT2022 dataset.

\subsection{Prediction and Detection Accuracy}

Table~\ref{tab:quant_ml} reports the accuracy metrics for each embedded model. Evaluation was performed using an 80/20 train--test split; anomaly detection metrics were computed against the 5\% injected DDoS-style outliers as ground-truth positives.

\begin{table}[H]
\centering
\caption{Model Accuracy Metrics (Test Set, $n = 209$)}
\label{tab:quant_ml}
\begin{tabular}{@{}llr@{}}
\toprule
\textbf{Model} & \textbf{Metric} & \textbf{Value} \\
\midrule
\multirow{3}{*}{Latency Predictor (Random Forest)}
    & MAE & 3.12 ms \\
    & RMSE & 4.67 ms \\
    & $R^2$ & 0.917 \\
\midrule
\multirow{3}{*}{Anomaly Detector (OneClassSVM)}
    & Precision & 0.873 \\
    & Recall & 0.841 \\
    & F1-Score & 0.857 \\
\midrule
\multirow{2}{*}{Forecaster (BiLSTM, t+10)}
    & MAE & 4.89 ms \\
    & $R^2$ & 0.862 \\
\bottomrule
\end{tabular}
\end{table}

The latency predictor achieves an $R^2$ of 0.917, indicating that the Random Forest captures most of the variance in flow duration from just two input features. The anomaly detector's F1 of 0.857 reflects a trade-off: precision is slightly higher than recall, meaning the model is conservative (it misses some injected outliers rather than over-flagging normal traffic). The forecaster's higher MAE compared to the direct predictor is expected given the 10-step horizon and compounding uncertainty.

\subsection{Federated Learning Convergence}

Table~\ref{tab:fed_convergence} tracks the global model MAE across 10 FedAvg communication rounds with $K=3$ edge nodes.

\begin{table}[H]
\centering
\caption{FedAvg Global Model Error per Communication Round}
\label{tab:fed_convergence}
\begin{tabular}{@{}crr@{}}
\toprule
\textbf{Round} & \textbf{Global MAE (ms)} & \textbf{$\Delta$ from Previous} \\
\midrule
1  & 11.43 & --- \\
2  & 9.87  & $-$1.56 \\
3  & 8.21  & $-$1.66 \\
4  & 7.14  & $-$1.07 \\
5  & 6.58  & $-$0.56 \\
6  & 6.19  & $-$0.39 \\
7  & 5.93  & $-$0.26 \\
8  & 5.81  & $-$0.12 \\
9  & 5.76  & $-$0.05 \\
10 & 5.74  & $-$0.02 \\
\bottomrule
\end{tabular}
\end{table}

The global MAE drops sharply in the first 4 rounds (from 11.43 to 7.14), then converges gradually. By round 8, marginal improvements fall below 0.15\,ms per round, suggesting that the current data partition and model capacity are approaching saturation. The final MAE of 5.74\,ms is higher than the centralized predictor's 3.12\,ms, reflecting the cost of distributing training across heterogeneous partitions without data sharing.

\subsection{Runtime Latency}

Table~\ref{tab:runtime} reports per-packet inference timing measured inside the AnyLogic JVM across 1,047 consecutive packets.

\begin{table}[H]
\centering
\caption{Per-Packet Inference Latency (JVM, $n = 1{,}047$)}
\label{tab:runtime}
\begin{tabular}{@{}lrrr@{}}
\toprule
\textbf{Component} & \textbf{Mean (\textmu s)} & \textbf{P95 (\textmu s)} & \textbf{P99 (\textmu s)} \\
\midrule
OfflineAiPredictor & 84 & 127 & 203 \\
AnomalyModel & 112 & 168 & 241 \\
FutureForecaster & 96 & 143 & 219 \\
QTableBalancer & 11 & 18 & 29 \\
\midrule
\textit{Total pipeline} & \textit{303} & \textit{456} & \textit{692} \\
\bottomrule
\end{tabular}
\end{table}

All models execute in under 1\,ms on average, with P99 latencies below 700\,\textmu s for the complete pipeline. The QTableBalancer is an order of magnitude faster than the ML models because it performs a simple hash-map lookup rather than tree traversal. These timings confirm that embedded inference adds negligible overhead relative to the simulation's event-processing cycle.

